\chapter*{Exit Review}
\markboth{Exit Review}{}
\addcontentsline{toc}{chapter}{Exit Review}

Solve these twelve problems by hand in one sitting, with no notes, then verify every numeric answer with your scene from Chapter~11 or the chapter functions it is built from. Eleven correct passes the course; fewer means rereading the chapters you missed and retaking it a week later. These answers are deliberately not in the key.

\begin{multicols}{2}
\begin{reviewexercises}
\item Write the contrapositive of ``If two lines are parallel, then corresponding angles are congruent,'' and state whether it is true.
\item Two parallel lines are cut by a transversal and one angle measures $52\degs$. Find its same-side interior partner and its alternate exterior partner.
\item Is a triangle with sides $10$, $24$, $26$ a right triangle? Find its area.
\item Find the distance and midpoint between $(-2, -3)$ and $(4, 5)$.
\item Rotate $(5, -2)$ by $90\degs$ counterclockwise about the origin, then reflect the result across the $y$-axis.
\item Find each interior angle and each exterior angle of a regular decagon.
\item Find the area of the polygon $(0, 0)$, $(6, 0)$, $(6, 4)$, $(2, 7)$ by the shoelace formula.
\item A circle has radius $9$. Find the arc length and sector area for a central angle of $80\degs$, exactly and to two decimals.
\item A 30-foot ladder leans against a wall at $72\degs$ to the ground. How high does it reach, to two decimals?
\item Convert $210\degs$ to radians and give $\sin$ and $\cos$ of that angle exactly.
\item Find the angle between $\vv{3, 4}$ and $\vv{-8, 6}$.
\item Find where the line $y = x$ meets the circle $x^2 + y^2 = 50$.
\end{reviewexercises}
\end{multicols}

\medskip
The curriculum's exit check for Course~2 also applies, and the review above covers it:

\begin{checklist}
\item Prove the Pythagorean theorem and the triangle angle sum.
\item Convert the standard angles between degrees and radians without a calculator.
\item Compute distance, midpoint and the dot product of two vectors, and explain what the sign of the dot product tells you.
\item Write a function that rotates a point about the origin by any angle.
\end{checklist}
